%%%%%
%
%% Dans ce fichier, les �ventuels accents sont cod�s
%% en ISO-8859-15 (aka latin9)
%% La ligne suivante donne trois exemples 
%% e accent aigu : �, e accent grave : �, e-dans-l'o : �
\section{R�ciproque de la vitesse conventionnelle a�ronautique}

\itemstyle{Op�rateur~:} \operfmt{aero\_inv\_vit\_conv}
\index{aero\_inv\_vit\_conv}

%%%%%OBJET%%%%
\itemstyle{Objet~:}
Calcul de la pression dynamique correspondant � une
vitesse conventionnelle donn�e.

%%%%%ARGUMENTS%%%%
\bigskip
%% Dans ce fichier, les �ventuels accents sont cod�s
%% en ISO-8859-15 (aka latin9)
%% La ligne suivante donne trois exemples 
%% e accent aigu : �, e accent grave : �, e-dans-l'o : �
%%%%%ENTREES%%%%
\begin{tabarg}
\multicolumn{3}{l}{\titretab {Entr\'ees}}\\
\hline
V$_c$ & Vecteur, \texttt{float} & Vitesse de l'air (kt ou m/s)\\
unit & Coeff., \texttt{String} & Unit� de V$_c$\\
\end{tabarg}

\medskip
%%%%%SORTIES%%%%
\begin{tabarg}
\multicolumn{3}{l}{\titretab {Sorties}}\\
\hline
$\Delta P$ & Vecteur, \texttt{float} & Pression dynamique correspondant 
              � V$_c$ dans le cadre des vitesses conventionnelles (hPa)\\
\end{tabarg}


%%%%%FORMULE%%%%
\itemstyle{Formule}

%% Formule de l'operateur : aero_inv_vit_conv
%% Dans ce fichier, les �ventuels accents sont cod�s
%% en ISO-8859-15 (aka latin9)
%% La ligne suivante donne trois exemples 
%% e accent aigu : �, e accent grave : �, e-dans-l'o : �
Si unit = "kt", alors $V_{s0} = 661.479$. Sinon, si
unit = "m/s", alors  $V_{s0} = 340.294 $
\[
\Delta P = 1013.25 \left[ \left(1+\frac{1}{5}\frac{V_c^2}{V_{s0}^2}\right)^{\frac{7}{2}} -1 \right]
\]



%%%%%AUTEUR%%%%
\itemstyle{Auteur~:}
CNRM/GMEI/TRAMM

%%%%%REFERENCE%%%%
\itemstyle{R�f�rences~:}

Norme AFNOR  NF X 02-115 : \emph{Symbole et vocabulaire de la m�canique
du vol} \cite{AFNOR-74}, paragraphe 5.6.1 page 45.

Cette norme fixe en particulier les valeurs num�riques de $P_0$ 
(1013.25 hPa), de la vitesse du son $a_n$ (340.294 m.s$^{-1}$),
et de $\gamma$ (1.400, ce qui induit $\frac{\gamma-1}{2} = 1/5$
et $\frac{\gamma}{\gamma-1} = 7/2$).
